\section{CI/CD Dasar: Otomasi Build dan Deploy (Konsep)}

\textbf{CI/CD} (\textit{Continuous Integration / Continuous Deployment}) adalah praktik mengotomasi proses build, test, dan deploy setiap kali ada perubahan kode. \textbf{CI} memastikan kode dari berbagai developer diintegrasikan dan diuji secara otomatis. \textbf{CD} mengotomasi deployment ke staging atau produksi setelah semua tes lulus. Alat populer: \textbf{GitHub Actions}, GitLab CI, Jenkins. CI/CD mengurangi human error, mempercepat rilis, dan memberikan kepercayaan bahwa kode yang di-deploy telah teruji \cite{mdnDeploy}.

\begin{figure}[ht]
\centering
\begin{tikzpicture}[
  box/.style={draw, rounded corners, minimum width=1.8cm, minimum height=0.8cm, align=center, font=\footnotesize, fill=#1},
  arrow/.style={-Stealth, thick},
  node distance=0.7cm
]
  \node[box=blue!15] (push) {Git Push};
  \node[box=green!15, right=0.7cm of push] (build) {Build};
  \node[box=orange!15, right=0.7cm of build] (test) {Test};
  \node[box=yellow!15, right=0.7cm of test] (stage) {Deploy\\Staging};
  \node[box=red!15, right=0.7cm of stage] (prod) {Deploy\\Prod};

  \draw[arrow] (push) -- (build);
  \draw[arrow] (build) -- (test);
  \draw[arrow] (test) -- (stage);
  \draw[arrow] (stage) -- (prod);

  \node[above=0.3cm of build, font=\footnotesize\bfseries, blue!60] {CI};
  \node[above=0.3cm of stage, font=\footnotesize\bfseries, red!60] {CD};
\end{tikzpicture}
\caption{Pipeline CI/CD: push $\to$ build $\to$ test $\to$ deploy}
\end{figure}

